\section{Скрипты для создания графики}

Чтобы изобразить множество \( P \) и все его промежуточные образы, воспользуемся графическим пакетом \textit{TikZ}. Для генерации необходимых файлов напишем скрипты на языке \textit{Python}.

Для генерации точек, описывающих множество \( P \), ограничим \( \abs{z}\leq 3 \). Ещё для наглядности разделим комплексную плоскость на 4 подмножества и окрасим их градиентами, которые меняются от близости к началу координат:

\begin{center}
  \begin{tblr}{
      colspec = {r Q[c,m] Q[c,m]},
      column{1} = {rightsep=6pt},
      column{2} = {leftsep=6pt},
      column{3} = {leftsep=12pt},
      hline{2-5} = {dashed},
      vline{2-3} = {dashed},
      row{1} = {font=\bfseries},
      row{2-Z} = {rowsep=4pt},
  }
      & \(\Arg z\) & Градиент \\
      1. & \(\left(-\pi, -\dfrac{3\pi}{4}\right)\) & 
      \begin{tikzpicture}[baseline=4pt]
          \shade[left color=red,right color=yellow] (0,0) rectangle (3,0.5);
      \end{tikzpicture} \\
      2. & \(\left[-\dfrac{3\pi}{4}, 0\right)\) &
      \begin{tikzpicture}[baseline=4pt]
          \shade[left color=black,right color=salmon] (0,0) rectangle (3,0.5);
      \end{tikzpicture} \\
      3. & \(\left[0, \dfrac{3\pi}{4}\right]\) &
      \begin{tikzpicture}[baseline=4pt]
          \shade[left color=blue,right color=green-cyan] (0,0) rectangle (3,0.5);
      \end{tikzpicture} \\
      4. & \(\left(\dfrac{3\pi}{4}, \pi\right)\) & 
      \begin{tikzpicture}[baseline=4pt]
          \shade[left color=cyan,right color=magenta] (0,0) rectangle (3,0.5);
      \end{tikzpicture} \\
  \end{tblr}
\end{center}

\lstinputlisting[
  language=Python,
  caption={Скрипт \texttt{generate\_dots.py}},
]{script/generate_dots.py}

Для транспиляции окрашенных наборов точек в графики \textit{TikZ} воспользуемся скриптом:

\lstinputlisting[
  language=Python,
  caption={Скрипт \texttt{tikzify.py}},
]{script/tikzify.py}